\documentclass[12pt]{article} 
\renewcommand{\thempfootnote}{\arabic{mpfootnote}}

\usepackage[hidelinks]{hyperref}

\usepackage[utf8]{inputenc}

\usepackage[T1]{fontenc}
\usepackage[ttdefault=true]{AnonymousPro}

\usepackage{multicol}
\setlength{\columnsep}{1cm}

\usepackage{abc}
\usepackage{graphicx}
\usepackage{musicography}

\usepackage[margin=0.5in]{geometry} 
\geometry{letterpaper}

\usepackage{fancyhdr} 
\pagestyle{fancy} % options: empty , plain , fancy
\renewcommand{\headrulewidth}{0pt} % customise the layout...
\lhead{}\chead{}\rhead{}
\lfoot{}\cfoot{\thepage}\rfoot{}

\begin{document}

\thispagestyle{empty}

\begin{abc}[name=morpethrant,options={-O= +c -F single.fmt}]
X:1
T:The Morpeth Rant
M:C|
K:D
|: (3ABc | "D"d2 AG FDFA | "G"BGBd "A"cAce | "D"f2 df gfed | "A"cdec A2 Bc | 
"D"d2 AG FDFA | "G"BGBd "A"cAce | "D"f2 df "G"gfed | "A"Ac "D"d2 d2 :|
|: (3ABc | "D"dfaf dfaf | "G"gfef g2 ef | "A"gfed cdeg | "D"fefg f2 e2 | 
"D"dfaf dfaf | "G"gfef g2 ef | "A"gfed cAag | "D"f2 d2 d2 :| 
\end{abc}

\vskip 0.5in
\noindent\makebox[\textwidth][c]{%
  \begin{minipage}{7.0in}
    \setlength{\parindent}{1.5em}

    The tune is attributed to William Shield (1748- 1829) from
    Swalwell, Northumberland, northern England, but in the great
    tradition of folk music, he probably was not the composer as it is
    now believed that he appropriated his
    tunes.\footnote{\url{https://tunearch.org/wiki/Annotation:Morpeth\_Rant\_(1)}}
    Morpeth is a market town also in Northumberland, so we can safely
    assume the tune is from that general area. The Morpeth Rant is
    sometimes called the Morpeth Hornpipe, a hornpipe being in 4/4,
    played a little slower than a reel, with the last bar having a
    rhythm of 2 quarter notes then 1 half note (or 3 quarter notes and
    a pickup into the next phrase). The Morpeth Rant is also the name
    of a dance that uses the Morpeth Rant tune or a variant as
    accompaniment, one version being used as a Morris Dance
    tune.\footnote{\textit{Loc. cit.}}  A good source for details
    regarding The Morpeth Rant dance and rant step can be found on
    Andrew Carnie’s folk music site,
    folkdancemusings.blogspot.com.\footnote{I quickly learned that
      there are many (passionately opposing) opinions regarding the
      rant step so I chose to sidestep the argument and direct you
      here:
      \url{https://folkdancemusings.blogspot.com/2020/01/morpeth-rant-england.html}}
    The Morpeth Rant can be heard in Old Time Jams, Irish and Scottish
    Sessions, and Contradances.

    ---Kitty

  \end{minipage}
  \vspace*{\fill}
}
  \vspace*{\fill}
\end{document}
